\exercisegroup

\begin{enumerate}
\def\labelenumi{\arabic{enumi})}
\item
  Let E, F and G be three locally convex spaces, S a cover of E consisting of bounded sets. Show that if u is a separately continuous bilinear mapping from \(E \times F\) into G such that for every set \(M \in S\) the restriction of u to \(M \times F\) is continuous, then u is S-hypocontinuous.
\item
  Let \(\mathbf{E}\) be the direct sum space \(\mathbf{R}^{(\mathbf{N})}\) , with the topology induced by the product topology on \(\mathbf{R}^{\mathbf{N}}\) . Show that the bilinear form \(((x_n), (y_n)) \mapsto \sum_{n=0}^{\infty} x_n y_n\) on \(\mathbf{E} \times \mathbf{E}\) , is separately continuous, but that for every set \(\mathfrak{S}\) of bounded subsets of \(\mathbf{E}\) containing at least one infinite dimensional bounded set this bilinear form is not \(\mathfrak{S}\) -hypocontinuous.
\end{enumerate}


\begin{enumerate}
\def\labelenumi{\arabic{enumi})}
\setcounter{enumi}{2}
\tightlist
\item
  Let E be the space \(\mathbf{R}^{(\mathbf{N})}\) with the finest locally convex topology (II, p.~26); let F be the space \(R^{N}\) ; the space E is ultrabornological (III, p.~45, exerc. 19) and complete, whilst F is metrizable and complete. Let S (resp. T) be the set of all bounded subsets of E (resp. F). Show that the bilinear form \(((x_{n}), (y_{n})) \mapsto \sum_{n=0}^{\infty} x_{n} y_{n}\) on E × F is \((\mathfrak{S}, \mathfrak{T})\) -hypocontinuous, but is not continuous
\end{enumerate}

(cf.~IV, p.~48, exerc. 11).

\begin{enumerate}
\def\labelenumi{\arabic{enumi})}
\setcounter{enumi}{3}
\item
  Let E be a locally convex space, F an infrabarrelled space (III, p.~44, exerc. 7) and T the set of all bounded subsets of F. Show that, if a bilinear mapping from \(E \times F\) into a locally convex space G is T-hypocontinuous, it is \((\mathfrak{S}, \mathfrak{T})\) -hypocontinuous for every set S of bounded subsets of E (cf.~III, p.~44, exerc. 11).
\item
  \begin{enumerate}
  \def\labelenumii{\alph{enumii})}
  \tightlist
  \item
    Let \(\mathbf{E}\) , \(\mathbf{F}\) and \(\mathbf{G}\) be three Hausdorff locally convex spaces, and \(u\) a bilinear mapping from \(\mathrm{E} \times \mathrm{F}\) into \(\mathbf{G}\) . In order that there exist a balanced neighbourhood \(\mathbf{U}\) of 0 in \(\mathbf{E}\) such that the set of all mappings \(u(x, .)\) , where \(x\) runs through \(\mathbf{U}\) , is equicontinuous in \(\mathcal{L}(\mathbf{F}; \mathbf{G})\) , it is necessary and sufficient that \(u\) is continuous when we replace the topology of \(\mathbf{E}\) by the coarsest topology for which the sets \(\lambda \mathbf{U} (\lambda \neq 0)\) form a fundamental system of neighbourhoods of 0. Show that if \(\mathbf{G}\) is normed, this condition is satisfied by every continuous bilinear mapping from \(\mathrm{E} \times \mathrm{F}\) into \(\mathbf{G}\) .
  \end{enumerate}
\end{enumerate}

\begin{enumerate}
\def\labelenumi{\alph{enumi})}
\setcounter{enumi}{1}
\tightlist
\item
  Take for E, F and G the product space \(R^{N}\) , and for u the continuous bilinear mapping \(((x_{n}), (y_{n})) \mapsto (x_{n}y_{n})\) . Show that there does not exist any neighbourhood U of 0 in E such that the set of maps \(u(x, .)\) , where x runs through U, is equicontinuous in \(\mathcal{L}(F; G)\) .
\end{enumerate}

\begin{enumerate}
\def\labelenumi{\arabic{enumi})}
\setcounter{enumi}{5}
\tightlist
\item
  Let E, F and G be three topological vector spaces. A set H of bilinear mappings from \(E \times F\) into G is said to be separately equicontinuous if for all \(x \in E\) , the set of linear mappings \(u(x, .)\) , where u runs through H, is equicontinuous in \(\mathcal{L}(F; G)\) and if for all \(y \in F\) , the set of linear mappings \(u(. , y)\) , where u runs through H, is equicontinuous in \(\mathcal{L}(E; G)\) .
\end{enumerate}

Suppose that F is metrizable, and that E is a Baire space (cf.~III, p.~43, exerc. 5 and V, p.~79, exerc. 15). Show that every separately equicontinuous set of bilinear mappings from \(E \times F\) into G is equicontinuous (cf.~III, p.~42, exerc. 11).

\begin{enumerate}
\def\labelenumi{\arabic{enumi})}
\setcounter{enumi}{6}
\tightlist
\item
  Let \(\mathrm{E}\) , \(\mathrm{F}\) and \(\mathrm{G}\) be three topological vector spaces, \(\mathfrak{S}\) a set of bounded subsets of \(\mathrm{E}\) , and \(\mathrm{H}\) a set of separately continuous bilinear mappings from \(\mathrm{E} \times \mathrm{F}\) into \(\mathrm{G}\) . The following properties are equivalent:
\end{enumerate}

\(\alpha)\) For every neighbourhood \(\mathbf{W}\) of 0 in \(\mathbf{G}\) and every set \(\mathbf{M} \in \mathfrak{S}\) , there exists a neighbourhood \(\mathbf{V}\) of 0 in \(\mathbf{F}\) such that \(u(\mathbf{M} \times \mathbf{V}) \subset \mathbf{W}\) for all \(u \in \mathbf{H}\) .

\(\beta)\) For every set \(\mathbf{M} \in \mathfrak{S}\) , the image of \(\mathbf{H} \times \mathbf{M}\) under the mapping \((u, x) \mapsto u(x, .)\) is an equicontinuous subset of \(\mathcal{L}(\mathrm{F}; \mathrm{G})\) .

\(\gamma)\) As \(u\) runs through H, the set of mappings \(y \mapsto u(., y)\) from F into \(\mathcal{L}_{\mathfrak{S}}(\mathrm{E}; \mathrm{G})\) is equicontinuous.

We then say that H is a S-equihypocontinuous set of bilinear mappings (separately continuous) from \(E \times F\) into G. Similarly for a set T of bounded subsets of F, we define the notions of a T-equihypocontinuous set and a \((\mathfrak{S}, \mathfrak{T})\) -equihypocontinuous set.

\begin{enumerate}
\def\labelenumi{\arabic{enumi})}
\setcounter{enumi}{7}
\item
  Let H be a S-equihypocontinuous set of bilinear mappings from \(E \times F\) into G (exerc. 7). For every subset \(M \in S\) , show that H is equicontinuous in \(M \times F\) ; moreover, for every bounded subset Q of F, the union of the sets \(u(\mathbf{M} \times \mathbf{Q})\) , where u runs through H, is bounded in G.
\item
  Let H be a \((\mathfrak{S}, \mathfrak{T})\) -equihypocontinuous set of bilinear mappings from \(E \times F\) into G (exerc. 7); show that for every pair of sets \(M \in S\) , \(N \in T\) , H is uniformly equicontinuous in \(M \times N\) .
\item
  Let \(\mathrm{E}_1\) , \(\mathrm{E}_2\) and \(\mathrm{F}\) be three topological vector spaces, \(\mathbf{G}_1\) (resp. \(\mathbf{G}_2\) ) an everywhere dense subspace of \(\mathrm{E}_1\) (resp. \(\mathrm{E}_2\) ), and \(\mathfrak{S}_1\) (resp. \(\mathfrak{S}_2\) ) a family of bounded subsets of \(\mathbf{G}_1\) (resp. \(\mathbf{G}_2\) ). Let \(\mathrm{H}\) be a set of separately continuous bilinear mappings from \(\mathrm{E}_1 \times \mathrm{E}_2\) into \(\mathrm{F}\) ; if the set of restrictions to \(\mathbf{G}_1 \times \mathbf{G}_2\) of the mappings \(u \in \mathrm{H}\) is \((\mathfrak{S}_1, \mathfrak{S}_2)\) -equihypocontinuous, then so is \(\mathrm{H}\) .
\item
  If \(\mathbf{F}\) is a barrelled space, every separately equicontinuous set of bilinear mappings from \(\mathbf{E} \times \mathbf{F}\) into a locally convex space \(\mathbf{G}\) is \(\mathfrak{S}\) -equihypocontinuous for every set \(\mathfrak{S}\) of bounded subsets of \(\mathbf{E}\) .
\item
  Let E, F be two topological vector spaces, and let f be the bilinear mapping \((x, u) \mapsto u(x)\) from \(E \times \mathcal{L}(E; F)\) into F; let T be a topology compatible with the vector space structure of \(\mathcal{L}(E; F)\) and finer than the topology of simple convergence. Let S be a family of bounded subsets of E, U a family of bounded subsets of \(\mathcal{L}(E; F)\) (for the topology T). Show that f is S-hypocontinuous if and only if T is finer than the S-topology; f is U-hypocontinuous if and only if the sets of U are equicontinuous subsets of \(\mathcal{L}(E; F)\) .
\item
  Let E, F, G be three topological vector spaces, and S (resp. T) a family of bounded subsets of E (resp. F). Let H be the vector space of T-hypocontinuous bilinear mappings from \(E \times F\) into G.
\end{enumerate}

\begin{enumerate}
\def\labelenumi{\alph{enumi})}
\item
  Show that on H the topology of uniform convergence on sets of the form \(M \times N\) , where \(M \in S\) and \(N \in T\) is compatible with the vector space structure; this topology is called the \((\mathfrak{S}, \mathfrak{T})\) -topology on H. For every mapping \(u \in H\) , let \(\tilde{u}\) be the continuous mapping \(x \mapsto u(x, .)\) from E into \(\mathcal{L}_{\mathfrak{T}}(F; G)\) . Show that \(u \mapsto \tilde{u}\) is an isomorphism from the space H, endowed with the \((\mathfrak{S}, \mathfrak{T})\) -topology, onto the space \(\mathcal{L}_{\mathfrak{S}}(E; \mathcal{L}_{\mathfrak{T}}(F; G))\) .
\item
  Let L be a subset of H such that, for every pair \((x, y) \in E \times F\) , the set of all \(u(x, y)\) , where u runs through L, is bounded in G (simply bounded subset of H). Show that, if E, F, G are locally convex, and if E and F are Hausdorff and quasi-complete, then L is bounded in H for the \((\mathfrak{S}, \mathfrak{T})\) -topology.
\item
  Let E, F, G be three Hausdorff locally convex spaces. If E is barrelled and F quasi-complete or barrelled, then every simply bounded subset L of H is T-equihypocontinuous (III, p.~47, exerc. 7).
\end{enumerate}

\(d\) ) If E and F are barrelled, and G quasi-complete, and if \(\mathfrak{S}\) and \(\mathfrak{T}\) are coverings of E and F respectively, then H is Hausdorff and quasi-complete for the \((\mathfrak{S},\mathfrak{T})\) -topology.

\begin{enumerate}
\def\labelenumi{\arabic{enumi})}
\setcounter{enumi}{13}
\item
  Extend the definitions and results of § 5 to arbitrary multilinear mappings. Let E, F, G be three topological vector spaces, S (resp. T) a family of bounded subsets of E (resp. F), and U a family of bounded subsets of the space \(\mathcal{L}_{\mathfrak{S},\mathfrak{T}}\) (E, F; G) of bilinear (S, T)-hypocontinuous mappings from E × F into G, endowed with the (S, T)-topology (exerc. 13). Show that the trilinear map (x, y, u) → u(x, y) from E × F × \(\mathcal{L}_{\mathfrak{S},\mathfrak{T}}\) (E, F; G) into G is (S, T)-hypocontinuous; in order that it is (S, U)-hypocontinuous, it is necessary and sufficient that every set L ∈ U is S-equihypocontinuous (III, p.~47, exerc. 7).
\item
  Let \(E\) be the space of all sequences \(x = (\xi_n)_{n \geqslant 0}\) of real numbers such that the series with the general term \(\xi_n\) is convergent. Put \(\| x \| = \sup_n | \sum_{k=0}^n \xi_k |\) .
\end{enumerate}

\begin{enumerate}
\def\labelenumi{\alph{enumi})}
\tightlist
\item
  Show that \(\| x\|\) is a norm on E, and that E is complete for this norm.
\end{enumerate}

\begin{enumerate}
\def\labelenumi{\alph{enumi})}
\setcounter{enumi}{1}
\tightlist
\item
  Show that the vector space \(\ell^1 (\mathbf{N})\) (I, p.~4), considered as a subspace of E, is everywhere dense (for the topology of E); the topology on \(\ell^1 (\mathbf{N})\) defined by the norm \(\| x\| _1 = \sum_{n = 0}^{\infty}|\xi_n|\) is strictly finer than the topology induced by that of E.
\end{enumerate}

\begin{enumerate}
\def\labelenumi{\alph{enumi})}
\setcounter{enumi}{2}
\tightlist
\item
  Let \((\mathbf{P}_n)\) be an increasing sequence of finite subsets of \(\mathbf{N} \times \mathbf{N}\) forming a cover of \(\mathbf{N} \times \mathbf{N}\) . For every \(x = (\xi_n) \in \mathrm{E}\) and every \(y = (\eta_n) \in \ell^1(\mathbf{N})\) , let \(f_n(x, y) = \sum_{(i,j) \in \mathbf{P}_n} \xi_i \eta_j\) . Then the sequence \((f_n(x, y))\) tends to a limit for every pair \((x, y) \in \mathrm{E} \times \ell^1(\mathbf{N})\) if and only if for each of these pairs \((x, y)\) , the sequence \((f_n(x, y))\) is bounded; the limit of \(f_n(x, y)\) is then equal to \(\left(\sum_{n=0}^{\infty} \xi_n\right) \left(\sum_{n=0}^{\infty} \eta_n\right)\) .
\end{enumerate}

(Using exerc. 13, c) of III, p.~48, show that the sequence of bilinear forms \((f_n)\) is equicontinuous, and observe that it converges in the subspace \(\ell^1 (\mathbf{N})\times \ell^1 (\mathbf{N})\) ; conclude the argument using \(b).\) \(d)\) For every \(j\in \mathbf{N}\) , let \(\rho_{jn}\) be the smallest number of closed intervals of \(\mathbf{N}\) whose union is the projection of \(\mathbf{P}_n\cap (\mathbf{N}^{\prime}\times \{j\})\) onto \(\mathbf{N}\) ; let \(\rho_{n} = \sup_{j\in \mathbf{N}}\rho_{jn}\) . Show that the condition obtained in \(c)\) is equivalent to \(\sup_n\rho_n < + \infty\) . (If \(\phi_n\) is the characteristic function of \(\mathbf{P}_n\) , show that

the norm of the bilinear form \(f_{n}\) is \(\sup_{j\in \mathbf{N}}\big(\sum_{i = 0}^{\infty}|\rho_n(i,j) - \phi_n(i + 1,j)|\big).\)

